% Part: history
% Chapter: biographies 
% Section: alan-turing 

\documentclass[../../../include/open-logic-section]{subfiles}

\begin{document}

\olfileid[pt]{his}{bio}{tur} 

\olsection{Alan Turing}

Alan Turing nasceu em Maida Vale, Londres, em 23 de junho de 1912. É
considerado o pai da ciência da computação teórica. O interesse de Turing
pelas ciências físicas e pela matemática começou em tenra idade. No entanto,
quando menino, seus interesses não eram bem contemplados nas escolas que
frequentava, onde a ênfase recaía sobre a literatura e os
clássicos. Consequentemente, ele teve um desempenho fraco na escola e foi
repreendido por muitos de seus professores.

\olphoto{turing-alan}{Alan Turing}

Turing estudou no King's College, em Cambridge, na graduação, onde
estudou matemática. Em 1936, Turing desenvolveu (o que hoje se chama)
a máquina de Turing como uma tentativa de definir precisamente a noção de
função computável e de provar a indecidibilidade do problema da decisão. Ele
foi superado no resultado por Alonzo Church, que o provou por meio de seu
próprio cálculo lambda. Ainda assim, o artigo de Turing foi publicado
com referência ao resultado de Church. Church convidou Turing para Princeton,
onde ele passou de 1936 a 1938 e obteve um doutorado sob a orientação de
Church.

Apesar de seu interesse pela lógica, os interesses anteriores de Turing pelas
ciências físicas permaneceram predominantes. Suas habilidades práticas foram
postas em prática durante seu serviço no departamento britânico de criptanálise
em Bletchley Park durante a Segunda Guerra Mundial. Turing foi uma figura
central na quebra da cifra usada nas comunicações navais alemãs---o código
Enigma. A expertise de Turing em estatística e criptografia, juntamente
com a introdução de maquinário eletrônico, deu à equipe a
capacidade de quebrar o código por meio da criação de uma máquina de
descriptografia chamada ``bombe''. Suas ideias também ajudaram na criação do
primeiro computador eletrônico programável do mundo, o Colossus, também usado
em Bletchley Park para quebrar a cifra alemã Lorenz.

Turing era homossexual. Mesmo assim, em 1942 ele pediu Joan Clarke em
casamento, uma de suas colegas de equipe em Bletchley Park, mas depois rompeu
o noivado e confessou a ela que era homossexual. Teve vários amantes
ao longo da vida, embora os atos homossexuais fossem, então, crimes no Reino
Unido. Em 1952, a casa de Turing foi assaltada por um amigo de seu amante da
época e, ao prestar queixa à polícia, Turing admitiu ter um relacionamento
homossexual, sob a impressão de que o governo estava a caminho de
legalizar os atos homossexuais. Isso não era verdade, e ele foi acusado de
grave indecência. Em vez de ir para a prisão, Turing optou por um tratamento
hormonal que reduzia a libido. Turing foi encontrado morto em 8 de junho de
1954, de uma overdose de cianeto---muito provavelmente suicídio. Ele recebeu
um perdão real da Rainha Elizabeth~II em 2013.

\begin{reading}
Para uma biografia abrangente de Alan Turing, veja \citet{Hodges2014}.
A vida e a obra de Turing inspiraram uma peça, \emph{Breaking the Code},
que foi produzida em 1996 para a TV, com Derek Jacobi no papel de
Turing. \emph{The Imitation Game}, um filme indicado ao Oscar,
com Bendict Cumberbatch e Kiera Knightley, também é baseado livremente
na vida de Alan Turing e em seu período em Bletchley Park
\citep{Imitation2014}.

\citet{Radiolab2012} tem vários podcasts sobre a vida e a obra de Turing.
O documentário \emph{The Strange Life and Death of Dr.~Turing}, da BBC
Horizon, está disponível para assistir on-line \citep{Sykes1992}.
\citep{Theelen2012} é um vídeo curto de uma Máquina de Turing de LEGO em
funcionamento---feito para homenagear o centenário de Turing em 2012.

O artigo original de Turing sobre máquinas de Turing e o problema da decisão é
\citet{Turing1937}.
\end{reading}

\end{document}
